\documentclass[10pt,a4paper]{amsart}
\usepackage[margin=19mm]{geometry}
\usepackage{amsmath,amssymb,mathtools}
\usepackage[scr=rsfs,cal=euler]{mathalfa}
\usepackage{xfrac}
\usepackage[unicode,hidelinks,bookmarks=false]{hyperref}
\newcommand{\ie}{i.e.\ }
\newcommand{\cf}{cf.\ }
\newcommand{\resp}{respectively\ }
\newcommand{\Subsection}{\subsection}
\providecommand{\nobreakdash}{\nobreak\hbox{-}}
\setlength{\emergencystretch}{2em}
\multlinegap0pt
\usepackage{amssymb}
\usepackage{xcolor}
\usepackage{tikz}
\usepackage{dsfont}
\usepackage[shortlabels]{enumitem}
\newtheorem{prop}{Proposition}
\title{Di-graphs with tightly connected clusters: Effective Graph Laplacians and resolvent convergence}
\author{Christian Koke}
\date{Source: arXiv:2601.18057v2 (CC BY 4.0)}
\begin{document}
\maketitle

\noindent\emph{Source and licence.} Di-graphs with tightly connected clusters: Effective Graph Laplacians and resolvent convergence. Version arXiv:2601.18057v2 (CC BY 4.0), distributed by arXiv under the Creative Commons Attribution 4.0 International licence, http://creativecommons.org/licenses/by/4.0/, as recorded in the arXiv licence metadata for this exact version and shown on the abstract page https://arxiv.org/abs/2601.18057v2. Full licence notice: \texttt{licenses/arxiv-2601.18057-CC-BY-4.0.txt}.\par\smallskip

\noindent\textit{Editorial scope.} Counted unit: the article's prop prop (source0.tex lines 604--609). The article's complete original proof of it is delivered with it (source0.tex lines 610--614, one \texttt{proof} environment of 5 lines). No mathematics is written, repaired or re-derived: the statement and the proof are the article's own lines, in the article's own notation, and no retained line is substituted or re-flowed.  Retained uncounted as the source-local context the proof needs: lines 589--591.  Nothing was omitted from any retained span, so the package carries no omission record.  No retained line contains a citation, so the wrapper carries no bibliography and no bibliography entry.  No retained line contains a figure environment, an image-inclusion command or drawing code, so no image is delivered and the package declares no asset dependency.  Editorial formatting only, in addition: an AMS \texttt{amsart} preamble that restates the article's own declarations and macros used by the retained lines, namely the environments \texttt{prop} on separate counters, together with the standard packages those lines need.  The retained environments print in this document as Proposition 1 in order of appearance, whereas the article prints the same items with the numbering of its own full document.  The complete native source of the article as distributed in the arXiv e-print for this exact version is bundled unchanged as \texttt{sources/193365/source0.tex}.
	\begin{align}\label{dirichlet_form}
		\mathcal{E}_G(f,f) :=  \sum\limits_{(i,j) \in E} a(i,j) |f(i) - f(j)|^2 = \langle f, L f\rangle_{\ell^2(G)}.
	\end{align}

\begin{prop}
	The kernel of $L$ on $G = (V, m, E, a)$ is given as 
	\begin{align}
	\ker(L) = \overline{\operatorname{span} \left\{ \chi_C : C \in \mathfrak{C} \text{ and } \underline{m}(C) < \infty \right\}}.
	\end{align}
\end{prop}

\begin{proof}
Clearly $\operatorname{span} \left\{ \chi_C : C \in \mathfrak{C} \text{ and } m(C) < \infty \right\} \subseteq \ker(L)$ and $\ker(L)$ is closed, so that "$\supseteq
$"  is proved.
To also prove "$\subseteq$", let now $f \in \ker(L)$ be arbitrary.  Since $f \in \ker(L)$,  clearly also $\langle f, L f \rangle_{\ell^2(G)} = 0$, and hence the Dirichlet form (\ref{dirichlet_form}) vanishes on $f$. Thus we see, that a necessary condition for $f$ to be in the kernel of $L$ is, that $f$ is constant on each connected component of $G$, as otherwise $\langle f, L f \rangle_{\ell^2(G)} = \mathcal{E}_G(f,f) > 0$. Suppose now there is a connected component $C$ of $G$ on which the restriction $f{\restriction_C} = \lambda \neq 0$ does not vanish. Then $|\lambda|^2 \cdot m(C) \leq \|f\|^2_{\ell^2(G)} < \infty$ and hence the mass of this connected component is finite.
\end{proof}

\end{document}
